% Order statistics of an iid sample: the densities of the minimum,
% a central value, and the maximum pull apart as n grows. Drawn for
% n = 5 from a Unif(0,1) parent so the order-statistic densities are
% Beta(k, n-k+1), which TikZ can plot in closed form.
\begin{figure}[htbp]
\centering
\begin{tikzpicture}[x=10cm, y=1.6cm, font=\small]

% Axes
\draw[->, thick] (-0.02, 0) -- (1.05, 0) node[right] {$x$};
\draw[->, thick] (0, -0.05) -- (0, 2.9) node[above] {density};

% Parent uniform density on [0,1] (height 1)
\draw[thick, enggray, dashed] (0, 1) -- (1, 1);
\node[anchor=west, enggray, font=\scriptsize] at (1.0, 1.0)
  {parent $\mathrm{Unif}(0,1)$};

% Beta(k, n-k+1) density for n=5: f(x) = [n!/((k-1)!(n-k)!)] x^{k-1}(1-x)^{n-k}
% Minimum: k=1 -> 5 (1-x)^4
\draw[thick, engblue, smooth]
  plot[domain=0:1, samples=80] (\x, {5*pow(1-\x,4)});
\node[anchor=south, engblue, font=\scriptsize] at (0.12, 2.55)
  {$X_{(1)}$ (min)};

% Median: k=3 -> 30 x^2 (1-x)^2
\draw[thick, engorange, smooth]
  plot[domain=0:1, samples=80] (\x, {30*pow(\x,2)*pow(1-\x,2)});
\node[anchor=south, engorange, font=\scriptsize] at (0.5, 1.95)
  {$X_{(3)}$ (median)};

% Maximum: k=5 -> 5 x^4
\draw[thick, engred, smooth]
  plot[domain=0:1, samples=80] (\x, {5*pow(\x,4)});
\node[anchor=south, engred, font=\scriptsize] at (0.88, 2.55)
  {$X_{(5)}$ (max)};

% x-axis ticks
\foreach \t in {0, 0.5, 1.0} {
  \draw (\t, 0.03) -- (\t, -0.03) node[anchor=north, font=\scriptsize] {\t};
}

\end{tikzpicture}
\caption[Order-statistic densities for a five-point uniform
sample]{Densities of the order statistics of an iid sample of size
$n = 5$ drawn from $\mathrm{Unif}(0, 1)$. Sorting the sample turns the
minimum, the median, and the maximum into distinct random variables;
for the uniform parent the $k$th order statistic is
$\mathrm{Beta}(k, n - k + 1)$. The minimum concentrates near $0$ and
the maximum near $1$, each with mean $1/(n+1)$ and $n/(n+1)$, while
the central order statistic is symmetric about $1/2$. The minimum is
the weakest-link variable of the chain example in the text; the spread
between maximum and minimum is the sample range used by range-based
control charts. Source: closed-form Beta densities, schematic by the
authors.}
\label{fig:vol02:ch13:order-statistics}
\end{figure}
